\documentclass{amsart}
\usepackage{amsmath,amssymb,amsthm,mathtools}
\usepackage{hyperref}
\usepackage[margin=1in]{geometry}
\usepackage{array}
\usepackage{booktabs}

\newtheorem{theorem}{Theorem}[section]
\newtheorem{conjecture}[theorem]{Conjecture}
\newtheorem{proposition}[theorem]{Proposition}
\newtheorem{corollary}[theorem]{Corollary}
\newtheorem{lemma}[theorem]{Lemma}
\theoremstyle{definition}
\newtheorem{definition}[theorem]{Definition}
\newtheorem{example}[theorem]{Example}
\theoremstyle{remark}
\newtheorem{remark}[theorem]{Remark}

\DeclareMathOperator{\sch}{\mathfrak{S}}
\DeclareMathOperator{\groth}{\mathfrak{G}}
\newcommand{\code}{\mathfrak{c}}
\newcommand{\tom}[1]{\ensuremath{\xrightarrow{#1}}}
\newcommand{\Tom}[1]{\ensuremath{\xRightarrow{#1}}}
\newcommand{\fgl}{\oplus}
\newcommand{\fgn}{\ominus}
\newcommand{\Fix}{\mathrm{Fix}}
\newcommand{\Lft}{\mathrm{Left}}
\newcommand{\Out}{\mathrm{Out}}
\newcommand{\Supp}{\mathrm{Supp}}
\newcommand{\ZZ}{\mathbb{Z}}
\newcommand{\schubmult}{\textsf{schubmult}}
\newcommand{\pyid}[1]{\texttt{#1}}

\title[An equivariant $K$-theoretic Molev--Sagan formula]{An equivariant $K$-theoretic Molev--Sagan type formula for double Grothendieck polynomials}
\author{Matthew J. Samuel}
\date{\today}

\begin{document}

\begin{abstract}
We describe the multiplication rule for double Grothendieck polynomials implemented in
\pyid{schubmult.mult.groth\_double}. It rests on a conjectural closed formula for the
coefficients of $\prod_{i\le k}(x_i+z)\,\groth_u(x;y)$ in the $\groth$-basis: each coefficient
is a factorial elementary symmetric polynomial $E_{n,n}(\mathcal{Y};-z)$ in an alphabet
$\mathcal{Y}$ built from the formal inverses $\fgn y_a$ of the fixed window values and one copy
of $-1/\beta$ per left-moving value, times a prefactor of powers of $\beta$ and atoms
$(1+\beta y_a)^{-1}$. At $\beta=0$ this is the Pieri formula of
\cite[Theorem~7.1]{samuelmolev}. Because the coefficient is itself a factorial elementary
symmetric polynomial, its $z$-divided differences are again of that form, which gives a
closed Pieri rule for every $E_{p,k}(x;z)$ and, through the elementary-symmetric expansion of
double Schubert polynomials, a Molev--Sagan type expansion of $\groth_u(x;y)\,\groth_v(x;z)$
in the basis $\{\groth_w(x;y)\}$. The formula has been verified exhaustively on $S_4$ and
sampled through $S_6$.
\end{abstract}

\maketitle

\section{Conventions}

Throughout, $\beta$ is a formal parameter and we write
\[
a \fgl b = a + b + \beta ab, \qquad \fgn a = \frac{-a}{1+\beta a},
\]
for the multiplicative formal group law and its formal inverse, so that
$a \fgl (\fgn a) = 0$ and $1+\beta(a\fgl b) = (1+\beta a)(1+\beta b)$. Two identities are used
repeatedly:
\begin{equation}\label{eq:fgl-ids}
    1 + \beta(\fgn a) = \frac{1}{1+\beta a},
    \qquad
    z \fgl (\fgn a) = \frac{z - a}{1+\beta a}.
\end{equation}

Permutations are written in one-line notation and $\ell(w)$ is the length (number of
inversions). For $u\le w$ in Bruhat order we write $\ell(u,w) = \ell(w)-\ell(u)$. The
transposition exchanging positions $a<b$ is $t_{ab}$.

\subsection*{Double Schubert polynomials}
$\sch_w(x;y)$ is the double Schubert polynomial in the convention of \cite{samuelmolev}: for
$w_0\in S_n$ the longest element,
$\sch_{w_0}(x;y) = \prod_{i+j\le n}(x_i - y_j)$, and $\sch_{ws_i} = \partial_i^x \sch_w$ when
$\ell(ws_i)<\ell(w)$. Thus $\sch_w(x;y)$ vanishes at $x_i = y_{w(i)}$ for all $i$, i.e.\ at the
torus fixed point indexed by $w$. This is the convention of \schubmult's \pyid{DSx}.

\subsection*{Double Grothendieck polynomials}
$\groth_w(x;y)$ is the double $\beta$-Grothendieck polynomial \cite{fomin1994grothendieck}
normalized so that
\[
\groth_{w_0}(x;y) = \prod_{i+j\le n}(x_i \fgl y_j),
\qquad
\groth_{ws_i} = \pi_i^x\,\groth_w \ \text{ if } \ell(ws_i)<\ell(w),
\]
where $\pi_i^x f = \partial_i^x\bigl((1+\beta x_{i+1}) f\bigr)$ is the isobaric divided
difference. This is the convention of \schubmult's \pyid{DGx}: for example
\[
\groth_{21}(x;y) = x_1\fgl y_1,\qquad
\groth_{312}(x;y) = (x_1\fgl y_1)(x_1\fgl y_2),\qquad
\groth_{231}(x;y) = (x_1\fgl y_1)(x_2\fgl y_1).
\]
With this normalization $\groth_w(x;y)$ vanishes at the point $x_i = \fgn y_{w(i)}$, and the
lowest-degree part of $\groth_w(x;y)$ is $\sch_w(x;-y)$. In particular
\begin{equation}\label{eq:beta0}
\groth_w(x;y)\big|_{\beta=0} = \sch_w(x;-y),
\end{equation}
so a formula for $\groth$ specializes to one for $\sch$ after $y\mapsto -y$.

\subsection*{Factorial elementary symmetric polynomials}
For $0\le p\le k$, $E_{p,k}(x;z)$ is the factorial elementary symmetric polynomial
\[
E_{p,k}(x;z) = \sum_{1\le i_1<\dots<i_p\le k}\ \prod_{j=1}^{p}\bigl(x_{i_j} - z_{i_j-j+1}\bigr)
= \sch_{c[p,k]}(x;z),
\qquad c[p,k] = s_{k-p+1}\cdots s_{k-1}s_k,
\]
so $E_{0,k}=1$ and $E_{k,k}(x;z) = \prod_{i=1}^{k}(x_i - z_1)$. In the variable $t=z_1$ the top
one is a generating function for the ordinary elementary symmetric polynomials,
\begin{equation}\label{eq:top-gen}
E_{k,k}(x;t) = \prod_{i=1}^{k}(x_i - t) = \sum_{p=0}^{k} (-t)^{k-p}\,e_p(x_1,\dots,x_k),
\end{equation}
and the others are obtained from it by divided differences in $z$:
\begin{equation}\label{eq:E-from-top}
E_{p,k}(x;z) = (-1)^{k-p}\,\partial^{z}_{k-p}\cdots\partial^{z}_{2}\partial^{z}_{1}\,E_{k,k}(x;z_1),
\qquad
\partial^z_{j} f = \frac{f - s^z_j f}{z_j - z_{j+1}}.
\end{equation}
On monomials in $t=z_1$ the operator $\partial^z_{q}\cdots\partial^z_1$ acts by
$t^m \mapsto h_{m-q}(z_1,\dots,z_{q+1})$, the complete homogeneous symmetric polynomial
(zero when $m<q$).

\section{The classical Pieri formula}

Sottile's $k$-Bruhat relation $u\tom{k} w$ holds when there are transpositions
$t_{a_1b_1},\dots,t_{a_pb_p}$ with $a_i\le k<b_i$, the lower indices $a_i$ pairwise distinct,
the upper indices $b_1\ge b_2\ge\cdots\ge b_p$, each step a Bruhat cover, and
$w = u t_{a_1b_1}\cdots t_{a_pb_p}$; then $\ell(u,w)=p$. Write
\[
P_k(u,w) = \{\,u(i)\ :\ i\le k,\ u(i)=w(i)\,\}
\]
for the values in the window $\{1,\dots,k\}$ that are left in place.

\begin{theorem}[{\cite[Theorem~7.1]{samuelmolev}}]\label{thm:pieri}
For $u\in S_\infty$, $k\ge1$ and $0\le p\le k$,
\[
\sch_u(x;y)\,E_{p,k}(x;z) \;=\; \sum_{u\tom{k} w} E_{p-\ell(u,w),\,k-\ell(u,w)}\bigl(y_{P_k(u,w)};z\bigr)\,\sch_w(x;y),
\]
where $y_{P_k(u,w)}$ denotes the alphabet $\{y_a : a\in P_k(u,w)\}$ ($E_{p,k}$ is symmetric in its
first alphabet, so the order is immaterial).
\end{theorem}

The case $p=k$ is a product formula: $E_{k-\ell,k-\ell}(y_P;z) = \prod_{a\in P}(y_a - z_1)$, i.e.
\begin{equation}\label{eq:pieri-top}
\sch_u(x;y)\prod_{i=1}^{k}(x_i - z) = \sum_{u\tom{k} w}\ \prod_{a\in P_k(u,w)}(y_a - z)\ \sch_w(x;y).
\end{equation}
Each factor has a localization meaning: at a fixed window position $i$ the polynomial $\sch_u$
``knows'' $x_i = y_{u(i)}$, and the factor $(x_i - z)$ contributes its value there. Moved
positions contribute $1$, and $u\tom{k}w$ forces exactly one unit of length per moved position.

\section{The $K$-theoretic top-block formula}

Fix $u$, $k$, and $w$ and look at the $k$ window positions. Each $i\le k$ falls into exactly
one of three classes according to the fate of the value $u(i)$ in $w$:
\begin{align*}
\Fix_k(u,w) &= \{\,i\le k : w(i) = u(i)\,\},\\
\Lft_k(u,w) &= \{\,i\le k : u(i) = w(j)\text{ for some } j<i,\ j\le k\,\}
\quad\text{(moves left inside the window)},\\
\Out_k(u,w) &= \{\,i\le k\,\}\setminus(\Fix_k\cup\Lft_k)
\quad\text{(leaves the window, or moves right inside it)}.
\end{align*}
Let $m_k(u,w) = k - |\Fix_k(u,w)|$ be the number of moved window positions and
$d = \ell(u,w)$. Attach to the pair $(u,w)$ the \emph{prefactor}
\begin{equation}\label{eq:pref}
\rho_k(u,w) \;=\; \beta^{\,\ell(u,w) - m_k(u,w)}\ (-\beta)^{|\Lft_k(u,w)|}\ \prod_{i\in\Out_k(u,w)}\frac{1}{1+\beta y_{u(i)}},
\end{equation}
and the alphabet of size $n_k(u,w) = |\Fix_k(u,w)| + |\Lft_k(u,w)| = k - |\Out_k(u,w)|$
\begin{equation}\label{eq:alphabet}
\mathcal{Y}_k(u,w) \;=\; \bigl(\fgn y_{u(i)}\bigr)_{i\in\Fix_k(u,w)}\ \cup\ \Bigl(-\tfrac{1}{\beta}\Bigr)^{|\Lft_k(u,w)|},
\end{equation}
the fixed values run through the formal inverse and each left-mover contributes the constant
$-1/\beta$.

\begin{definition}[The support]\label{def:supp}
$\Supp_k(u)$ is the set of $w$ such that $\groth_w(x)$ occurs in
$\prod_{i\le k}(x_i+z)\,\groth_u(x)$, the product of \emph{single} Grothendieck polynomials
with $z$ a formal variable. Equivalently, $\Supp_k(u) = \bigcup_{p=0}^{k}\Supp\bigl(e_p(x_1,\dots,x_k)\groth_u(x)\bigr)$,
which is the set of endpoints of the marked $k$-Pieri chains of Lenart--Sottile
\cite{lenart2007pieri}.
\end{definition}

\begin{conjecture}[Equivariant $K$-theoretic Pieri formula, top block]\label{conj:top}
For every $u\in S_\infty$, $k\ge1$ and formal $z$,
\begin{equation}\label{eq:top}
\prod_{i=1}^{k}(x_i + z)\ \groth_u(x;y)
\;=\;
\sum_{w\in\Supp_k(u)} \rho_k(u,w)\ E_{n,n}\bigl(\mathcal{Y}_k(u,w);\,-z\bigr)\ \groth_w(x;y),
\qquad n = n_k(u,w).
\end{equation}
\end{conjecture}

Since $E_{n,n}(\mathcal{Y};-z) = \prod_{Y\in\mathcal{Y}}(Y+z)$, and
$(-\beta)\bigl(-\tfrac1\beta + z\bigr) = 1-\beta z$, the summand of \eqref{eq:top} is the product
\begin{equation}\label{eq:top-product}
\beta^{\,\ell(u,w) - m_k(u,w)}
\prod_{i\in\Fix_k(u,w)}\bigl(z + \fgn y_{u(i)}\bigr)
\prod_{i\in\Lft_k(u,w)}\bigl(1 - \beta z\bigr)
\prod_{i\in\Out_k(u,w)}\frac{1}{1+\beta y_{u(i)}},
\end{equation}
which is the form in which the rule is implemented. Using \eqref{eq:fgl-ids} the three kinds of
factors can be written uniformly as
\[
z + \fgn y_{u(i)} = \frac{z(1+\beta y_{u(i)}) - y_{u(i)}}{1+\beta y_{u(i)}},
\qquad
1 - \beta z = 1 + \beta(-z),
\qquad
\frac{1}{1+\beta y_{u(i)}} = 1 + \beta(\fgn y_{u(i)}).
\]
In particular every coefficient is a polynomial in $\beta, z, y$ divided by a product of the
atoms $1+\beta y_a$, one for each moved-or-fixed window value except the left-movers; the
$\beta^{-1}$ in the alphabet is always cancelled by the $(-\beta)^{|\Lft|}$ in $\rho_k$.

\begin{remark}[Why $-1/\beta$ and not $1$]
One is tempted to write the left-mover factor $1-\beta z$ as ``$1 - (\beta z)$'' and use the
alphabet $(\fgn y_{\Fix}, 1^{|\Lft|})$ against a second alphabet mixing $-z$ and $\beta z$. That
product is not a factorial elementary symmetric polynomial in a single $z$-alphabet, and the
divided differences of Section~\ref{sec:Epk} would then have to be worked out term by term.
Rescaling instead by $-\beta$ puts every factor in the shape $Y - (-z)$ with a common second
variable, which is exactly what makes $E_{n,n}(\mathcal{Y};-z)$ available and, with it, the
closed form of Corollary~\ref{cor:Epk}.
\end{remark}

\begin{remark}[Consistency checks]
\begin{enumerate}
\item \emph{Specialization $\beta=0$.} In the product form \eqref{eq:top-product} the powers $\beta^{\ell(u,w)-m_k(u,w)}$ kill every $w$ with
$\ell(u,w) > m_k(u,w)$; what survives are the $w$ reached by exactly one Bruhat cover per moved
window position, which are Sottile's chains $u\tom{k}w$. The fixed factors become $z - y_{u(i)}$ and the others $1$. Substituting
$y\mapsto -y$ and $z\mapsto -z$ and using \eqref{eq:beta0} gives \eqref{eq:pieri-top}, i.e.\
Theorem~\ref{thm:pieri} with $p=k$.
\item \emph{The diagonal term.} Evaluating both sides of \eqref{eq:top} at the fixed point
$x_i=\fgn y_{u(i)}$ kills every $\groth_w$ with $w\ne u$ on the right (only $w\le u$ survive
localization at $u$, and $w\in\Supp_k(u)$ forces $w\ge u$). The left side becomes
$\prod_{i\le k}(z+\fgn y_{u(i)})\,\groth_u(\fgn y_u;y)$, which is the $w=u$ term of
\eqref{eq:top}. So the coefficient of $\groth_u$ is correct unconditionally; it is the
$K$-theoretic form of the localization reading of \eqref{eq:pieri-top}.
\item \emph{Non-equivariant limit.} At $y=0$ the fixed factors are $z$, the out-factors are
$1$, and \eqref{eq:top} reads
$\prod(x_i+z)\groth_u = \sum_w \beta^{\ell(u,w)-m}\,z^{|\Fix|}(1-\beta z)^{|\Lft|}\groth_w$.
Expanding in $z$ recovers the non-equivariant Pieri coefficients of
\cite{lenart2007pieri}, and shows that the equivariant and non-equivariant supports coincide,
which is why $\Supp_k(u)$ could be defined non-equivariantly.
\end{enumerate}
\end{remark}

\begin{remark}[On the support]
The support is genuinely smaller than the set of endpoints of Bruhat chains
$u = w_0 \lessdot w_1 \lessdot \cdots \lessdot w_d = w$ by transpositions $t_{a_jb_j}$ with
$a_j\le k<b_j$ and weakly decreasing $b_j$ (lower indices may now repeat). For instance with
$u=1243$ and $k=2$ the permutation $w=3412$ is such an endpoint, but
$w\notin\Supp_2(1243)$; the product formula would assign it the nonzero value
$\beta/((1+\beta y_1)(1+\beta y_2))$, so the restriction to $\Supp_k(u)$ in \eqref{eq:top}
cannot be dropped. The implementation enumerates the marked chains of \cite{lenart2007pieri}
(\pyid{schub\_lib.elem\_sym\_chains\_groth}) to decide membership.
\end{remark}

\begin{example}\label{ex:top21}
$u=21$, $k=2$. The window values are $u(1)=2$, $u(2)=1$, and $\Supp_2(21)=\{21,231,312,321\}$:
\[
\begin{array}{c c c c l}
\toprule
w & \Fix & \Out & \ell(u,w)-m & \text{coefficient of } \groth_w(x;y) \\
\midrule
21  & \{1,2\} & \emptyset & 0 & (z+\fgn y_2)(z+\fgn y_1) \\
231 & \{1\}   & \{2\}     & 0 & (z+\fgn y_2)\,\dfrac{1}{1+\beta y_1} \\
312 & \{2\}   & \{1\}     & 0 & (z+\fgn y_1)\,\dfrac{1}{1+\beta y_2} \\
321 & \emptyset & \{1,2\} & 0 & \dfrac{1}{(1+\beta y_1)(1+\beta y_2)} \\
\bottomrule
\end{array}
\]
Hence
\[
(x_1+z)(x_2+z)\,\groth_{21}(x;y)
= (z+\fgn y_1)(z+\fgn y_2)\,\groth_{21}
+ \frac{z+\fgn y_2}{1+\beta y_1}\,\groth_{231}
+ \frac{z+\fgn y_1}{1+\beta y_2}\,\groth_{312}
+ \frac{\groth_{321}}{(1+\beta y_1)(1+\beta y_2)}.
\]
\end{example}

\begin{example}\label{ex:top1243}
$u=1243$, $k=2$ shows all three classes and a nontrivial $\beta$-power. Here $u(1)=1$,
$u(2)=2$, $\ell(u)=1$, and $\Supp_2(1243)$ has seven elements:
\[
\begin{array}{c c c c c l}
\toprule
w & \Fix & \Lft & \Out & \ell(u,w)-m & \text{coefficient} \\
\midrule
1243 & \{1,2\} & & & 0 & (z+\fgn y_1)(z+\fgn y_2) \\
1342 & \{1\} & & \{2\} & 0 & (z+\fgn y_1)/(1+\beta y_2) \\
1423 & \{1\} & & \{2\} & 0 & (z+\fgn y_1)/(1+\beta y_2) \\
1432 & \{1\} & & \{2\} & 1 & \beta\,(z+\fgn y_1)/(1+\beta y_2) \\
2341 & & \{2\} & \{1\} & 0 & (1-\beta z)/(1+\beta y_1) \\
2413 & & \{2\} & \{1\} & 0 & (1-\beta z)/(1+\beta y_1) \\
2431 & & \{2\} & \{1\} & 1 & \beta\,(1-\beta z)/(1+\beta y_1) \\
\bottomrule
\end{array}
\]
In the last three rows the value $2$ has moved from position $2$ to position $1$, hence is a
left-mover; the value $1$ has left the window.
\end{example}

\section{From the top block to every $E_{p,k}$}\label{sec:Epk}

Conjecture~\ref{conj:top} is stated for a single formal variable $z$, but that is enough to
determine the coefficients of every $E_{p,k}(x;z)\groth_u$, exactly as in the Molev--Sagan
approach to factorial Schur functions \cite{molev1999littlewood}: by \eqref{eq:E-from-top},
$E_{p,k}(x;z)$ is obtained from $E_{k,k}(x;z_1)$ by divided differences in $z$, and those
operators can be pushed onto the coefficients in \eqref{eq:top}. The point of writing the
coefficient as $\rho_k(u,w)\,E_{n,n}(\mathcal{Y};-z)$ is that its $z$-divided differences are
again factorial elementary symmetric polynomials, by the very same identity
\eqref{eq:E-from-top} applied to the alphabet $\mathcal{Y}$ instead of $x$.

\begin{corollary}[Pieri rule for $E_{p,k}$]\label{cor:Epk}
Assume Conjecture~\ref{conj:top}. For $0\le p\le k$, $q=k-p$, and any alphabet $z_1,z_2,\dots$,
\[
E_{p,k}(x;z)\,\groth_u(x;y) \;=\; \sum_{w\in\Supp_k(u)} \kappa_k(u,w;q;z)\ \groth_w(x;y),
\qquad
\kappa_k(u,w;q;z) \;=\; \rho_k(u,w)\ E_{\,n-q,\,n}\bigl(\mathcal{Y}_k(u,w);\ z_1,\dots,z_{q+1}\bigr),
\]
with $n=n_k(u,w)$ and $\rho_k$, $\mathcal{Y}_k$ as in \eqref{eq:pref}, \eqref{eq:alphabet}
($E_{n-q,n}=0$ when $q>n$).
\end{corollary}

\begin{proof}
Put $t=z_1$ in \eqref{eq:top} (so $-z=-t$): $E_{k,k}(x;t)\,\groth_u = \sum_w \rho_k(u,w)\,E_{n,n}(\mathcal{Y};t)\,\groth_w$,
an identity of polynomials in $t$ with coefficients in the field generated by $y$ and $\beta$
(the $\rho_k$ and $\mathcal{Y}$ do not involve $t$). Apply the operator
$(-1)^{q}\partial^z_{q}\cdots\partial^z_1$ to both sides. On the left it produces
$E_{p,k}(x;z)\groth_u$ by \eqref{eq:E-from-top}; on the right it acts only on
$E_{n,n}(\mathcal{Y};z_1)$, where \eqref{eq:E-from-top} with $x$ replaced by $\mathcal{Y}$
gives $E_{n-q,n}(\mathcal{Y};z_1,\dots,z_{q+1})$. Linear independence of the $\groth_w(x;y)$
finishes the proof.
\end{proof}

Setting every $z_j=0$ gives the ordinary elementary symmetric polynomials, since
$E_{r,n}(\mathcal{Y};0)=e_r(\mathcal{Y})$.

\begin{corollary}[Pieri rule for $e_p$]\label{cor:ep}
Assume Conjecture~\ref{conj:top}. For $0\le p\le k$,
\[
e_p(x_1,\dots,x_k)\,\groth_u(x;y) \;=\; \sum_{w\in\Supp_k(u)} \rho_k(u,w)\ e_{\,p-|\Out_k(u,w)|}\bigl(\mathcal{Y}_k(u,w)\bigr)\ \groth_w(x;y).
\]
\end{corollary}

\begin{remark}[Specialization to Theorem~\ref{thm:pieri}]
At $\beta=0$ only $w$ with $\ell(u,w)=m_k(u,w)$ survive in $\rho_k$. Write $L=|\Lft_k(u,w)|$
and $P=P_k(u,w)$, so $n = |P|+L$ and $|P| = k-m$. In
$(-\beta)^{L}E_{n-q,n}(\fgn y_P,(-1/\beta)^{L};z)$ the only terms that survive as $\beta\to0$ are
those using all $L$ copies of $-1/\beta$, and each such copy contributes
$(-\beta)(-1/\beta - z_j) = 1+\beta z_j\to1$; what remains is $E_{n-q-L,\,n-L}(-y_P;z)
= E_{p-m,\,k-m}(-y_P;z)$. Since $u\tom{k}w$ has $\ell(u,w)=m$, after $y\mapsto-y$ this is
exactly the coefficient $E_{p-\ell,k-\ell}(y_P;z)$ of Theorem~\ref{thm:pieri}.
Corollary~\ref{cor:Epk} is thus the $K$-theoretic form of Theorem~\ref{thm:pieri}: the
alphabet $y_P$ of fixed values becomes $\fgn y_P$ padded by one $-1/\beta$ per left-mover, and
the prefactor $\rho_k$ carries the excess-length power of $\beta$ and one atom
$(1+\beta y_a)^{-1}$ per value that leaves the window.
\end{remark}

\begin{example}
Continuing Example~\ref{ex:top21} ($u=21$, $k=2$; no left-movers, so every alphabet consists of
formal inverses):
\[
\begin{array}{c c c c}
\toprule
w & \rho_2(21,w) & \mathcal{Y}_2(21,w) & n \\
\midrule
21  & 1 & (\fgn y_2,\fgn y_1) & 2\\
231 & (1+\beta y_1)^{-1} & (\fgn y_2) & 1\\
312 & (1+\beta y_2)^{-1} & (\fgn y_1) & 1\\
321 & \bigl((1+\beta y_1)(1+\beta y_2)\bigr)^{-1} & () & 0\\
\bottomrule
\end{array}
\]
Corollary~\ref{cor:ep} with $p=1$ reads $e_1(\fgn y_1,\fgn y_2)\groth_{21} + e_0\,\groth_{231}/(1+\beta y_1) + e_0\,\groth_{312}/(1+\beta y_2)$
($321$ has $p-|\Out| = -1$ and drops out):
\[
e_1(x_1,x_2)\,\groth_{21}(x;y) = (\fgn y_1 + \fgn y_2)\,\groth_{21} + \frac{\groth_{231}}{1+\beta y_1} + \frac{\groth_{312}}{1+\beta y_2},
\]
while $p=2$ gives $e_2(x_1,x_2)\groth_{21} = (\fgn y_1)(\fgn y_2)\groth_{21} + \frac{\fgn y_2}{1+\beta y_1}\groth_{231}
+ \frac{\fgn y_1}{1+\beta y_2}\groth_{312} + \frac{\groth_{321}}{(1+\beta y_1)(1+\beta y_2)}$.
For $E_{1,2}(x;z) = x_1+x_2-z_1-z_2$ ($q=1$), Corollary~\ref{cor:Epk} replaces
$e_1(\fgn y_1,\fgn y_2)$ by $E_{1,2}(\fgn y_1,\fgn y_2;z_1,z_2) = \fgn y_1+\fgn y_2 - z_1 - z_2$
and leaves the other two coefficients ($E_{0,1}=1$) unchanged:
\[
\groth_{21}(x;y)\,E_{1,2}(x;z) = \bigl(\fgn y_1 + \fgn y_2 - z_1 - z_2\bigr)\groth_{21}
+ \frac{\groth_{231}}{1+\beta y_1} + \frac{\groth_{312}}{1+\beta y_2}.
\]
A left-mover appears in Example~\ref{ex:top1243}: for $w=2341$ one has $\rho_2 = (-\beta)/(1+\beta y_1)$,
$\mathcal{Y}_2 = (-1/\beta)$, $n=1$, so the $E_{2,2}(x;z)$ coefficient is
$(-\beta)(-1/\beta - z_1)/(1+\beta y_1) = (1+\beta z_1)/(1+\beta y_1)$ and the $E_{1,2}(x;z)$
coefficient is $(-\beta)E_{0,1}/(1+\beta y_1) = -\beta/(1+\beta y_1)$.
\end{example}

\section{The Molev--Sagan type formula}

The last ingredient is the expansion of a double Schubert polynomial into products of
factorial elementary symmetric polynomials, which is what drives the double Schubert
algorithm in \schubmult{} \cite{samuelmolev}. For $v\in S_\infty$ let $\theta(v)$ be the
code of the dominant permutation $\mathrm{dom}(v)$ obtained by repeatedly sorting the code of
$v$ (Definition~\ref{def:theta} below), and write $\theta_i = \theta_i(v^{-1})$. For
permutations $v',v''$ and $i\ge1$ write $v'\Tom{i}v''$ if $v'' = v't_{i,b_1}\cdots t_{i,b_r}$
with distinct $b_j\ne i$ and each step a Bruhat cover; set $\sigma_i(v',v'') = \#\{j : b_j<i\}$
and $D_i(v',v'') = \{v'(i)\}\cup\{v'(j) : v'(j)\ne v''(j)\}$.

\begin{definition}\label{def:theta}
$\mu$ is dominant if $\code_i(\mu)\ge\code_{i+1}(\mu)$ for all $i$. If $v$ is not dominant let
$i$ be maximal with $\code_i(v)<\code_{i+1}(v)$ and set $\mathrm{dom}(v)=\mathrm{dom}(vs_i)$;
otherwise $\mathrm{dom}(v)=v$. Then $\theta_i(v) = \code_i(\mathrm{dom}(v))$.
\end{definition}

\begin{theorem}[{\cite{samuelmolev}}]\label{thm:vpath}
For $v\in S_\infty$, with $r$ the number of nonzero parts of $\theta(v^{-1})$,
\[
\sch_v(x;z) = \sum_{(v_0,\dots,v_r)}\ \prod_{i=1}^{r}(-1)^{\sigma_i(v_{i-1},v_i)}\,
E_{\theta_i - \ell(v_{i-1},v_i),\ \theta_i}\bigl(x;\ z_{D_i(v_{i-1},v_i)}\bigr),
\]
the sum over chains $1 = v_0\Tom{1}v_1\Tom{2}\cdots\Tom{r}v_r = v\,\mathrm{dom}(v^{-1})$, where
$z_D$ is the alphabet $\{z_a : a\in D\}$ in the order produced by \pyid{call\_zvars}.
\end{theorem}

Each factor is an $E_{p,k}(x;\cdot)$ with $k=\theta_i$ and $p = \theta_i - \ell(v_{i-1},v_i)$, so
Corollary~\ref{cor:Epk} can be applied layer by layer.

\begin{theorem}[$K$-theoretic Molev--Sagan formula, conditional on Conjecture~\ref{conj:top}]\label{thm:ms}
For $u,v\in S_\infty$,
\[
\groth_u(x;y)\,\sch_v(x;z) = \sum_{w} c^{w}_{u,v}(y;z)\ \groth_w(x;y),
\qquad
c^{w}_{u,v}(y;z) = \sum_{\substack{(v_0,\dots,v_r)\\ (u_0,\dots,u_r)}}\ \prod_{i=1}^{r}
(-1)^{\sigma_i(v_{i-1},v_i)}\,
\kappa_{\theta_i}\bigl(u_{i-1},u_i;\ \ell(v_{i-1},v_i);\ z_{D_i(v_{i-1},v_i)}\bigr),
\]
where the $v$-chain is as in Theorem~\ref{thm:vpath}, the $u$-chain has $u_0=u$, $u_r=w$ and
$u_i\in\Supp_{\theta_i}(u_{i-1})$, and $\kappa$ is the closed-form coefficient of
Corollary~\ref{cor:Epk},
\[
\kappa_k(u',u'';q;z) = \rho_k(u',u'')\ E_{\,n_k(u',u'')-q,\ n_k(u',u'')}\bigl(\mathcal{Y}_k(u',u'');\ z_1,\dots,z_{q+1}\bigr).
\]
Consequently, if $\groth_v(x;z) = \sum_{v'} a_{v,v'}(z;\beta)\,\sch_{v'}(x;z)$ is the (finite)
expansion of $\groth_v$ into double Schubert polynomials, then
\[
\groth_u(x;y)\,\groth_v(x;z) = \sum_{w}\Bigl(\sum_{v'} a_{v,v'}(z;\beta)\,c^{w}_{u,v'}(y;z)\Bigr)\groth_w(x;y).
\]
\end{theorem}

\begin{proof}[Proof, given Conjecture~\ref{conj:top}]
Substitute Theorem~\ref{thm:vpath} for $\sch_v(x;z)$ and expand $\groth_u\cdot\prod_i E_i$
one factor at a time from $i=1$ to $r$ using Corollary~\ref{cor:Epk}, which is a formal
consequence of \eqref{eq:top} and \eqref{eq:E-from-top}. Each application replaces
$\groth_{u_{i-1}}$ by $\sum_{u_i}\kappa_{\theta_i}(u_{i-1},u_i;\ldots)\groth_{u_i}$. The second
statement is bilinearity.
\end{proof}

At $\beta=0$ every $\kappa$ becomes the factorial elementary symmetric polynomial in the fixed
values, the $u$-chains become chains of Sottile relations $u_{i-1}\tom{\theta_i}u_i$, and
Theorem~\ref{thm:ms} is the Molev--Sagan type formula of \cite{samuelmolev} (after
$y\mapsto-y$). The $K$-theoretic novelty is entirely in $\kappa$, and it is of the same shape
as the classical coefficient: a factorial elementary symmetric polynomial, now in the alphabet
$\fgn y_{\Fix}$ padded with one $-1/\beta$ per left-mover, times the prefactor $\rho$ carrying
the excess-length power of $\beta$ and the atoms $(1+\beta y_a)^{-1}$ for values that exit or
move right; and the $u$-steps are restricted to the marked-chain supports.

\begin{example}[Headline example]\label{ex:full}
$u=21$, $v=132$. The expansion of $\groth_{132}(x;z)$ into double Schubert polynomials is
\[
\groth_{132}(x;z) = (1+\beta z_1)^2(1+\beta z_2)\,\sch_{132}(x;z) + \beta(1+\beta z_1)(1+\beta z_2)\,\sch_{231}(x;z)
+ (z_1\fgl z_2)\bigl(2 + \beta(z_1\fgl z_2)\bigr),
\]
and $\theta(132^{-1}) = \theta(231^{-1}) = (2)$, so both $\sch_{132}(x;z) = E_{1,2}(x;z)$ and
$\sch_{231}(x;z) = E_{2,2}(x;z)$ are single layers with $k=2$. Assembling them with the coefficient
polynomials of Example~\ref{ex:top21} and simplifying,
\[
\groth_{21}(x;y)\,\groth_{132}(x;z) \;=\; \gamma\,\groth_{21}(x;y) + (1+\beta\gamma)\bigl(\groth_{231}(x;y)+\groth_{312}(x;y)+\beta\,\groth_{321}(x;y)\bigr),
\]
\[
\gamma \;=\; z_1\fgl z_2\fgl(\fgn y_1)\fgl(\fgn y_2) \;=\; \frac{(z_1\fgl z_2) - (y_1\fgl y_2)}{(1+\beta y_1)(1+\beta y_2)},
\qquad
1+\beta\gamma = \frac{(1+\beta z_1)(1+\beta z_2)}{(1+\beta y_1)(1+\beta y_2)}.
\]
At $\beta=0$ this is $\sch_{21}(x;-y)\sch_{132}(x;z) = (z_1+z_2-y_1-y_2)\sch_{21}(x;-y) + \sch_{231}(x;-y)+\sch_{312}(x;-y)$,
the familiar double Schubert product. The coefficient of $\groth_{21}$ is the formal-group-law
sum of the same four terms, and the factor $1+\beta\gamma$ multiplying the three higher terms is
the ``$K$-theoretic $1$''.
\end{example}

\section{Implementation and verification}

In \pyid{schubmult.mult.groth\_double}:
\begin{itemize}
\item \pyid{grothmult\_double\_top} implements \eqref{eq:top} (\pyid{\_top\_block\_coeff} is the
product of factors, \pyid{\_top\_block\_support} is $\Supp_k(u)$);
\item \pyid{\_groth\_elem\_sym\_frac} implements $\kappa$ of Theorem~\ref{thm:ms} in the closed form of
Corollary~\ref{cor:Epk}: it evaluates $E_{n-q,n}(\mathcal{Y};z)$ by the subset expansion
$\sum_{i_1<\dots<i_p}\prod_j(Y_{i_j}-z_{i_j-j+1})$ as a dynamic program over the alphabet
\eqref{eq:alphabet}, clearing the denominator of each $\fgn y_a$ entry as it is used so that the
numerator stays polynomial;
\item \pyid{\_groth\_schub\_vpath\_mul} implements $c^w_{u,v}$ by running the $u$-chains against
the precomputed $v$-path dictionaries (\pyid{compute\_vpathdicts}), exactly as
\pyid{schubmult\_double} does for double Schubert polynomials;
\item \pyid{grothmult\_double} implements the final displayed formula of Theorem~\ref{thm:ms},
with \pyid{dgroth\_to\_dschub} supplying the coefficients $a_{v,v'}$.
\end{itemize}
Coefficients are carried throughout as ``flat fractions'' (a polynomial numerator together
with exponents of the atoms $1+\beta y_a$), which is possible precisely because
\eqref{eq:top} produces no other denominators. The ring interface is
\pyid{DGx(u) * DGx(v, z)}, and \pyid{.simplify()} returns the coefficients in the factored
form used above.

Conjecture~\ref{conj:top} was checked against the exact fold
\pyid{grothmult\_double\_block} (which multiplies by one linear factor $x_i+z$ at a time using the
$K_T$-Chevalley formula of \cite{lenart2007affine}) for every $u\in S_4$ and $1\le k\le 3$, and for
samples in $S_5$ and $S_6$ with $k\le5$; Theorem~\ref{thm:ms} was checked by expanding both
sides as polynomials for every pair $u,v\in S_3$ and for samples of pairs in $S_4$. All checks passed. A proof of
Conjecture~\ref{conj:top} is expected to follow from the localization characterization of
the coefficients: the diagonal term is already forced, and the factors for moved positions
should follow by induction on $\ell(u,w)$ from the $K_T$-Chevalley formula.

\bibliographystyle{acm}
\bibliography{formschub}

\end{document}
